\documentclass[11pt]{article}
\usepackage[margin=1in]{geometry}
\usepackage{amsmath,amssymb,microtype}
\usepackage[hidelinks]{hyperref}

\newcommand{\Riesz}{(\mathrm R_{p,\gamma})}
\newcommand{\Reverse}{(\mathrm {RR}_{p,\gamma})}

\title{The Vicsek-graph Conjecture C of arXiv:2502.10837\\
is answered by arXiv:2606.05475}
\author{Literature-status note}
\date{13 August 2026}

\begin{document}
\maketitle

\begin{center}
\fbox{\parbox{0.91\textwidth}{\textbf{Status.}
The later Theorem~1.3 proves source Conjecture~C in a stronger form and
disproves the companion extension proposed immediately after it.  It gives a
complete classification away from one critical exponent; equality remains
open.  This is a literature implication, not a new proof.}}
\end{center}

\section{The source question}

For a Vicsek graph with diffusion escape exponent $\beta$, arXiv:2502.10837v3
compares
\[
 \Riesz:\quad \|\nabla f\|_p\leq C\|\Delta^\gamma f\|_p,
 \qquad
 \Reverse:\quad \|\Delta^\gamma f\|_p\leq C\|\nabla f\|_p.
\]
Conjecture~C on PDF page~8 states that, for $p>2$, $\Reverse$ holds whenever
\begin{equation}\label{eq:theta}
 \gamma>\theta(p):=\frac1\beta+\frac1p\left(2-\frac2\beta\right).
\end{equation}
The following comment proposes the companion assertion that, for $1<p<2$,
$\Riesz$ holds whenever $\gamma<\theta(p)$, and calls equality open
\cite{Fen25}.

\section{The later classification and the exact implication}

For a $D$-dimensional Vicsek graph one has $\beta=D+1$.  Theorem~1.3 of
arXiv:2606.05475v2 (PDF page~5) sets
\begin{equation}\label{eq:star}
 \gamma_*(p)=\frac1{D+1}+\frac1p\left(1-\frac2{D+1}\right)
 =\frac1\beta+\frac1p\left(1-\frac2\beta\right)
\end{equation}
and proves, for every $1<p<\infty$ and away from equality,
\[
 \begin{array}{c|cc}
 &\Riesz&\Reverse\\ \hline
 0<\gamma<\gamma_*(p)&\text{holds}&\text{fails}\\
 \gamma_*(p)<\gamma<1&\text{fails}&\text{holds}.
 \end{array}
\]
The paper explicitly says that its main goal is to prove the Devyver--Russ
reverse-transform conjecture in the discrete Vicsek setting \cite{Fen26}.

The source and later thresholds satisfy the elementary identity
\[
 \theta(p)-\gamma_*(p)=\frac1p>0.
\]
Consequently $\gamma>\theta(p)$ implies $\gamma>\gamma_*(p)$, so the later
theorem proves source Conjecture~C, in fact on the strictly larger range
$\gamma>\gamma_*(p)$.

The proposed companion extension is false.  Since $\gamma_*(p)<1$, for every
$1<p<2$ one can choose
\[
 \gamma_*(p)<\gamma<\min\{\theta(p),1\}.
\]
This $\gamma$ satisfies the proposed hypothesis $\gamma<\theta(p)$, but
Theorem~1.3 says that $\Riesz$ fails.

\section{Remaining scope}

The precise remaining question is the single line $\gamma=\gamma_*(p)$, not
the larger threshold $\theta(p)$ printed in the 2025 proposal.  The 2026
paper expressly leaves that critical case open.  A separate lane packet for
arXiv:2606.05475 records new partial endpoint information; it does not claim
the endpoint classification.

Thus the defensible status is
\[
 \boxed{\text{source Conjecture C: yes; proposed companion: no;
 critical line $\gamma=\gamma_*(p)$: open.}}
\]

\begin{thebibliography}{9}
\bibitem{Fen25}
J.~Feneuil,
\emph{In spaces with a slow diffusion, the Riesz transform is unbounded on
$L^p$, $p\in(2,\infty)$}, arXiv:2502.10837v3 (2025), Conjecture~C.

\bibitem{Fen26}
J.~Feneuil,
\emph{Reverse inequalities for super-Riesz transforms on graphs with a slow
diffusion}, arXiv:2606.05475v2 (2026), Theorem~1.3.
\end{thebibliography}

\end{document}
